\documentclass[11pt,a4paper]{article}

\usepackage[T1]{fontenc}
\usepackage[utf8]{inputenc}
\usepackage[brazil]{babel}
\usepackage{lmodern}
\usepackage{microtype}
\usepackage[margin=2cm]{geometry}
\usepackage{hyperref}
\usepackage{titlesec}
\usepackage{xcolor}

\hypersetup{
  colorlinks=true,
  linkcolor=black,
  urlcolor=black,
  citecolor=black,
  pdfborder={0 0 0}
}

\definecolor{sectioncolor}{HTML}{111111}
\definecolor{accentcolor}{HTML}{1A1A1A}

\pagestyle{empty}
\setlength{\parindent}{0pt}
\setlength{\parskip}{5pt}
\linespread{1.05}
\selectfont

\titleformat{\section}
  {\bfseries\large\color{sectioncolor}}
  {}{0pt}{}[\vspace{-2pt}{\color{sectioncolor}\titlerule[0.8pt]}]
\titlespacing*{\section}{0pt}{12pt}{8pt}

\newcommand{\cvrole}[4]{%
  \vspace{4pt}%
  {\bfseries\color{accentcolor}#1}\hfill{\itshape #2}\\[-2pt]%
  {\itshape #3}\hfill #4\\[6pt]%
}

\begin{document}

\begin{center}
  {\fontsize{20}{24}\selectfont\bfseries Kamille Hartmann Maccagnan}\\[8pt]
  {\large Analista de BI \textbar{} Power BI \& SQL \textbar{} Indicadores e Dashboards}\\[10pt]
  Campo Comprido, Curitiba-PR \enspace\textbullet\enspace (41) 9 9928-4424 \enspace\textbullet\enspace
  \href{mailto:kamillehartmannm@gmail.com}{kamillehartmannm@gmail.com}\\[3pt]
  \href{https://linkedin.com/in/kamille-hartmann-maccagnan/}{linkedin.com/in/kamille-hartmann-maccagnan/}
\end{center}

\vspace{6pt}

\section{Resumo Profissional}

Profissional de Business Intelligence com experiência na coleta, organização e análise de dados, desenvolvimento de dashboards em Power BI e construção de indicadores de desempenho em ambientes operacionais e corporativos. Atua na extração e transformação de dados com Power Query e SQL, automação de rotinas analíticas e validação de integridade das informações para apoio à tomada de decisão. Experiência na modelagem de bases de dados, elaboração de relatórios gerenciais e monitoramento de KPIs operacionais e financeiros. Perfil analítico, organizado e colaborativo, com vivência em projetos multidisciplinares e interface entre áreas de negócio e tecnologia.

\section{Experiência Profissional}

\cvrole{Anjuss}{Analista de Suporte de TI}{07/2025 -- Atual}

Atuo na análise de dados e apoio operacional com desenvolvimento e manutenção de dashboards em Power BI para monitoramento de indicadores e desempenho de processos. Realizo extração, tratamento e transformação de dados com Power Query e Excel, aplicando regras de validação para garantir precisão e consistência das informações utilizadas nas análises. Implemento automações de coleta e processamento de dados com Power Automate integradas ao Power BI, otimizando o tempo de entrega de relatórios. Produzo relatórios gerenciais que sintetizam resultados e apoiam a identificação de oportunidades de melhoria em processos operacionais e na qualidade das informações.

\cvrole{HORSE do Brasil}{Estagiária de TI (IS/IT LATAM)}{07/2024 -- 07/2025}

Atuei com governança de dados e suporte analítico em ambiente multinacional, elaborando relatórios executivos e dashboards para acompanhamento de performance de projetos e iniciativas estratégicas. Participei de projetos multidisciplinares com foco em indicadores, documentação de processos, padronização de informações e análise de riscos. Utilizei o ServiceNow para gestão de demandas e apliquei o Power Automate em rotinas de automação de processos. Atuei na interface entre áreas de negócio, fornecedores e stakeholders, contribuindo com dados e análises para apoio à gestão e tomada de decisão.

\cvrole{Restaurante Familiar}{Analista de Dados e Suporte Técnico}{01/2023 -- 07/2024}

Atuei na coleta, organização e análise de dados de vendas, performance operacional e resultados financeiros para apoio à gestão do negócio. Estruturei do zero a modelagem da base de dados em Excel e implementei dashboards em Power BI para monitoramento de KPIs e acompanhamento de resultados. Realizei análises diagnósticas de desempenho, identifiquei desvios e oportunidades de melhoria em processos operacionais, elaborando relatórios com conclusões e recomendações para a liderança. Atuei na validação e padronização de fontes de dados para garantir integridade e confiabilidade das análises.

\section{Competências Técnicas}

\textbf{BI e Análise de Dados:} Power BI (DAX, Power Query), SQL, Excel avançado, Google Sheets, análise de dados, estatística descritiva, modelagem de dados, KPIs, indicadores operacionais e financeiros, relatórios gerenciais, dashboards, visualização de dados.

\textbf{ETL, Automação e Qualidade de Dados:} Extração, tratamento e transformação de dados, bancos de dados relacionais, validação e integridade de dados, padronização de informações, automação de processos analíticos (Power Automate), governança de dados.

\textbf{Colaboração e Projetos:} Interface entre áreas de negócio e tecnologia, projetos multidisciplinares, comunicação com lideranças, gestão de demandas, documentação de processos, melhoria contínua.

\textbf{Ferramentas:} ServiceNow, Trello, Monday, Spotfire, Pacote Office.

\textbf{Metodologias:} Scrum, Kanban.

\textbf{Idiomas:} Inglês Avançado \enspace\textbullet\enspace Espanhol Básico \enspace\textbullet\enspace Coreano Básico

\section{Projetos e Conquistas}

Estruturação de base de dados e dashboards em Power BI para monitoramento de KPIs de vendas, performance operacional e resultados financeiros. Automação e padronização de fluxos de coleta, tratamento e análise de dados com Power Automate integrado ao Power BI. 1º lugar no Transformation Day Renault 2023 com solução orientada a dados e melhoria de processos. Participação no Hackathon Bosch 2023 com foco em resolução de problemas de negócio com uso de dados. Desenvolvimento de relatórios executivos e indicadores para apoio à tomada de decisão.

\section{Formação Acadêmica}

\textbf{PUCPR} \hfill Curitiba, PR\\
Graduação em Gestão da Tecnologia da Informação \hfill Previsão de conclusão: 06/2026\\[6pt]

\textbf{Escola Conquer} \hfill Curitiba, PR\\
Pós-graduação em Liderança e Gestão em Tecnologia \hfill Conclusão: 2024\\[6pt]

\textbf{Universidade Tuiuti do Paraná} \hfill Curitiba, PR\\
Graduação em Medicina Veterinária \hfill 06/2019\\[6pt]

\textbf{Qualittas} \hfill Curitiba, PR\\
Pós-graduação em Clínica Médica e Cirúrgica de Animais Exóticos e Pets Não Convencionais \hfill 12/2022

\end{document}
